% Transcription of folder 147, batch 04 (archivists' pages 61-80).
% Pass: Opus 5.5 (claude-opus-5-5), 2026-10-09 — first pass, unchecked against the pages by a human.
%
% Grothendieck numbers every other sheet (author's pp. 30-39 on archivists'
% pp. 61, 63, 65, ...); the sheets in between continue the text.
\documentclass[11pt,a4paper]{article}
\input{../preamble/grothendieck}

\folder{147}
\batch{4}
\pages{61}{80}
\foldertitle{[Structure à l'infini des Mg,ν, pages 1 à 90, dont table des matières] : notes manuscrites (s.d.).}
\dating{s.d. — le groupe « Autour de La "Longue Marche" à travers la théorie de Galois » (141 à 149) est daté [à partir de 1978]-1983}
\watermark{Édition de démonstration}

\begin{document}

\page{61}\note{p. 30 de l'auteur. La phrase commence au feuillet précédent, hors de ce lot.}
la courbe modulaire $\widehat{\mathfrak{X}}^{!}_{0,4}$ sur $\widehat{M}^{!}_{0,4}$, munie de ses quatre sections $s_0, s_1, s_\infty, s_\lambda$ (ou $t_1, t_2, t_3, t_4$), et \ill{} l'opération de $\mathfrak{S}_4$ \struck{dessus} sur $\widehat{\mathfrak{X}}^{!}_{0,4} \to \widehat{M}^{!}_{0,4}$, échangeant ces quatre sections entre elles. Je vais pas le faire ici -- cela donnerait des \uncertain{compl.} les \struck{schémas} \add{multiplicités} modulaires $\widehat{M}_{0,4}$ et $\widehat{\mathfrak{X}}_{0,4}$. Pour $\widehat{M}_{0,4}$, il y a lieu de remplacer le diagramme (76) par un diagramme purement numérique
\[
(78)\qquad
\begin{array}{c|c}
\text{codim } 0 & \text{codim } 1 \\ \hline
\overset{4}{\bullet} \;\xrightarrow{\;1\;} & \overset{2}{\bullet}\!\text{---}\!\overset{2}{\bullet}
\end{array}
\]
il n'y a plus qu'un seul pt à l'infini, correspondant aux \uncertain{deux} pts à l'infini de $\widehat{M}^{!}_{0,4}$, permutés \uncertain{transitivement} par $\mathfrak{S}_4$.

Caractéristiques d'EP fractionnaires. On trouve
\begin{align*}
\chi_!(M^{!\,\mathrm{an}}_{0,4}) &= \chi_!(U^{\mathrm{an}}_{0,3}) = 2 - 2\underset{0}{g} - \underset{3}{\nu} = -1 \\
\chi_!(\widehat{M}^{!\,\mathrm{an}}_{0,4}) &= \chi_!(\mathbb{P}^{1\,\mathrm{an}}) = 2 \\
\chi_!(M^{\mathrm{an}}_{0,4}) &= \chi_!(\mathbb{B}^{\mathrm{an}}_{0,3}, \mathfrak{S}_4)
= \chi_!\bigl(U_{0,3} \setminus \{-j, -\bar{j}\} \setminus \{\tfrac12, 2, -1\}, \mathfrak{S}_4\bigr) \\
&\qquad + \chi_!(\{-j, -\bar{j}\}, \mathfrak{S}_4) + \chi_!(\{\tfrac12, 2, -1\}, \mathfrak{S}_4) \\
&= -1 + \tfrac13 + \tfrac12 = -\tfrac{1}{24}
\end{align*}
\note{dans la troisième ligne, un accent (chapeau) biffé sur le $M$ du membre de gauche ; à la dernière ligne, le dénominateur du résultat est surchargé, lu $24$ ; tel qu'écrit, $-1 + \frac13 + \frac12$ ne fait pas $-\frac1{24}$ : les termes sont sans doute à diviser par l'ordre de $\mathfrak{S}_4$, ce que la page ne dit pas}
\marginal{\emph{NB} On peut écrire directement $\chi_!(U^{\mathrm{an}}_{0,3}) : [\mathfrak{S}_4] = -\frac{1}{24}$}

\page{62}
\[
\chi_!(\widehat{M}^{\mathrm{an}}_{0,4}) = \underbrace{\chi_!(U^{\mathrm{an}}_{0,3}, \mathfrak{S}_4)}_{-1/24} + \underbrace{\chi_!(\{0, 1, \infty\}, \mathfrak{S}_4)}_{1/8} = \frac{1}{12} .
\]
De façon plus fine, dans un \og groupe de Grothendieck\fg{} convenable, on aura
\begin{align*}
[M^{!}_{0,4}] &= [U_{0,3}] = \underbrace{T(1)}_{\text{Tate} \simeq [\mathbb{E}^1]} - 2 \\
[M_{0,4}] &= \tfrac14\bigl(T(1) - 2\bigr) + \tfrac{1}{12} + \tfrac18 = \tfrac14\bigl(T(1) - 1\bigr) - \tfrac{1}{24} \\
[\widehat{M}_{0,4}] &= [M_{0,4}] + \tfrac18 = \tfrac14\bigl(T(1) - 1\bigr) + \tfrac{1}{12}
\end{align*}
\note{sous « Tate » : « (d'\uncertain{augmentation} 1) » ; dans les deuxième et troisième lignes, un premier essai de calcul est raturé avant le résultat, illisible ; dans la deuxième, le facteur $\frac14$ est écrit en surcharge}
\marginal{\uncertain{valable} au-dessus de $\mathbb{Z}[\frac16]$ \uncertain{79} -- \ill{} \ill{} faire le calcul \ill{} au-dessus de $\mathbb{Z}$…}
On trouve \ill{} ces trois quantités \uncertain{à des} dénominateurs \uncertain{près} : $-1, -\frac{1}{24}, \frac{1}{12}$.

\textbf{2°) \emph{Cas des} $\widehat{M}_{1,1}$.}
\[
(77)\qquad
\begin{array}{c|c}
\text{codim } 0 & \text{codim } 1 \\ \hline
\underset{1}{\overset{1}{\bullet}} \;\xrightarrow{\;1\;} & \text{(boucle)}\overset{1}{\bullet}
\end{array}
\]
\drawing{en codimension 1, une boucle fermée sur un sommet marqué 1}
Il y a exactement un point à l'infini.

On a un morphisme (\emph{au-dessus de} $\mathbb{Z}[\frac12]$)
\[
(78)\qquad \widehat{M}_{1,1} \longrightarrow \widehat{M}'_{1,4} \overset{\mathrm{def}}{=} (\widehat{M}^{!}_{1,4}, \mathfrak{S}_3)
\]
\note{les indices se lisent « $1,4$ » ; la suite ($[U_{0,3}, \mathfrak{S}_3]$, $M_{0,4}$) suggère $0,4$. Le numéro (78) est déjà employé p. 61}
qui fait de $\widehat{M}_{1,1}$ une gerbe (de groupe \uncertain{$\{\pm 1\}$}) au-dessus de $\widehat{M}'_{1,4}$ (\uncertain{en car.} $\neq 2$), d'où
\[
[M_{1,1}] = \tfrac12 [U_{0,3}, \mathfrak{S}_3] = \tfrac12 [M_{0,4}] \text{\struck{$\ldots$}}
\]
\note{la fin de la ligne, où l'on devine $[\mathbb{E}^1]$ et $T(1)$, est raturée}

\page{63}\note{p. 31 de l'auteur}
d'où, au-dessus de $\mathbb{Z}[\frac12]$
\[
(79)\quad
\begin{cases}
[M_{1,1}] = \tfrac12\bigl(T(1) - 1\bigr) - \tfrac{1}{12} \\
[\widehat{M}_{1,1}] = [M_{1,1}] + \underset{1/2}{[\widehat{M}^{\infty}_{1,1}]} = \tfrac12\bigl(T(1) - 1\bigr) + \tfrac{5}{12}
\end{cases}
\]
\note{dans les deux lignes, un premier essai raturé précède le résultat ; dans la seconde on y devine $T(1)$ et $\frac13$}

Voici un autre calcul de $[M_{1,1}]$. On considère le morphisme modulaire
\[
M_{1,1} \longrightarrow \mathbb{E}^1
\]
\struck{et} on découpe $\mathbb{E}^1$ suivant les deux sections $S_0$ (correspondant à la courbe elliptique \uncertain{harmonique}, de groupe d'automorphismes $\mathbb{Z}/4\mathbb{Z}$) et $S_1$ \struck{\ill{}} (courbe elliptique \uncertain{équianharmonique}, \ill{} groupe d'automorphismes $\mathbb{Z}/6\mathbb{Z}$), qui ne se coupent qu'en les pts $2$ et $3$, qui correspondent (\ill{}) à des groupes d'automorphismes d'ordre $24$, et $12$. On a donc une partition
\[
\mathbb{E}^1 = U \amalg S'_0 \amalg S'_1 \cup a_2 \cup a_3
\]
\note{un premier symbole raturé après le signe $=$}
où
\[
\begin{cases}
U = \mathbb{E}^1 \setminus (S_0 \cup S_1) \\
\{a_2, a_3\} = S_0 \cap S_1 \quad (a_2 \text{ en car. } 2,\ a_3 \text{ en car. } 3) \\
S'_0 = S_0 \setminus \{a_2, a_3\},\quad S'_1 = S_1 \setminus \{a_2, a_3\}
\end{cases}
\]
On trouve donc

\page{64}
\[
[M_{1,1}] = \tfrac12 [U] + \tfrac14 [S'_0] + \tfrac16 [S'_1] + \tfrac{1}{24} [a_2] + \tfrac{1}{12} [a_3]
\]
et on a
\begin{align*}
[U] &= [\mathbb{E}^1] - [S_0] - [S_1] + [a_2] + [a_3] \\
&= T(1) - 2 + \zeta_2 + \zeta_3
\end{align*}
où on a posé en général, pour $p$ premier,
\[
\zeta_p = [\operatorname{Spec} \mathbb{Z}/p\mathbb{Z}] = i_{p!}(1)
\]
où $i_{p!} : K_!(\mathbb{Z}_p) \longrightarrow K_!(\mathbb{Z})$.

De même
\begin{align*}
[S'_0] &= [S_0] - [a_2] - [a_3] = 1 - \zeta_2 - \zeta_3 \\
[S'_1] &= [S_1] - [a_2] - [a_3] = 1 - \zeta_2 - \zeta_3
\end{align*}
d'où
\begin{align*}
[M_{1,1}] &= \tfrac12\bigl(T(1) - 2 + \zeta_2 + \zeta_3\bigr) + \underbrace{\bigl(\tfrac14 + \tfrac16\bigr)}_{5/12}\bigl(1 - \zeta_2 - \zeta_3\bigr) \\
&\qquad + \tfrac{1}{24}\zeta_2 + \tfrac{1}{12}\zeta_3 \\
&= \tfrac12 T(1) - \tfrac{7}{12} + \tfrac18 \zeta_2 + \tfrac16 \zeta_3 \\
&= \tfrac12\bigl(T(1) - 1\bigr) - \tfrac{1}{12} + \tfrac18 \zeta_2 + \tfrac16 \zeta_3 ,
\end{align*}
\note{dans la parenthèse $(\frac14 + \frac16)$ et dans la deuxième ligne, des chiffres surchargés ($\frac18$ récrit)}
donc
\[
(80)\quad
\begin{cases}
[M_{1,1}] = \tfrac12\bigl(T(1) - 1\bigr) - \tfrac{1}{12} + \tfrac18 \zeta_2 + \tfrac16 \zeta_3 \\
[\widehat{M}_{1,1}] = [M_{1,1}] + \tfrac12 = \tfrac12 T(1) - \tfrac{1}{12} + \tfrac18 \zeta_2 + \tfrac16 \zeta_3
\end{cases}
\]

\page{65}\note{p. 32 de l'auteur}
\textbf{3°) \emph{Cas des} $\widehat{M}^{!}_{0,5}$, ou $\widehat{M}_{0,I}$.}
\[
(81)\qquad
\begin{array}{c|c|c}
\text{codim } 0 & \text{codim } 1 & \text{codim } 2 \\ \hline
\overset{I(5)}{\bullet} \;\xrightarrow{\;(10)\;} &
\overset{I'(2)}{\bullet}\!\text{---}\!\overset{I''(3)}{\bullet} \;\xrightarrow{\;(3)\;} &
\overset{I'(2)}{\bullet}\!\text{---}\!\overset{I'''(1)}{\bullet}\!\text{---}\!\overset{I''(2)}{\bullet}
\end{array}
\]
\note{les accents des deux derniers indices de la codimension 2 sont surchargés}
\emph{Dans} $\widehat{M}^{!}_{0,5}$ il y a dix \og diviseurs\fg{} à l'infini \uncertain{($\simeq (\{0,1,\infty\}, I'')$)}, \uncertain{dont chacun est} un $\widehat{M}^{!}_{0,4}$ \struck{\ill{}}, \uncertain{plus précisément}, un $\widehat{M}_{0, I''(3) \amalg (*)} \simeq (\mathbb{P}^1_{I''})_{\mathbb{Z}}$ \struck{$\widehat{M}_{0,3}$, \ill{}} (\ill{} \ill{} et \uncertain{ayant}) trois points à l'infini -- les \ill{} sont en tout $\frac{30}{2} = 15$, correspondant aux partitions de $I(5)$ de type $(2,2,1)$, \uncertain{alors} que les diviseurs \ill{} correspondent aux parties $I'$ de $I$ de cardinal $2$. Deux tels diviseurs (correspondant à des parties $I', I'' \in \mathfrak{P}_2(I)$) se rencontrent \struck{elles} ssi $I' \cap I'' = \emptyset$, i.e. ssi $I'$, $I''$ et $I \setminus (I' \cup I'')$ \struck{forment} forment une partition de $I$ de type $(2,2,1)$.
\marginal{$\mathbb{P}^1_{I''} \overset{\mathrm{def}}{=} \mathbb{P}^1_{\mathbb{Z}} \wedge \mathfrak{S}_3$ \ill{}}

D'après l'étude combinatoire de l'icosaèdre, \uncertain{choisissant} une orientation de $I$ -- ce qui \uncertain{détermine} un icosaèdre orienté, \ill{} un \uncertain{icos. grand}) les \add{10} diviseurs à l'infini de $\widehat{M}_{0,I}$, i.e. les éléments de $\mathfrak{P}_2(I)$, \add{ou \uncertain{(les}} \struck{\uncertain{sommets}}

\page{66}
\uncertain{pris grand complémentaire}) aux dix \og triangles\fg{} de $I$, correspondant aux dix faces de l'icosaèdre grand (${\simeq}$ paires de faces opposées de l'icos. ordinaire), tandis que les 15 partitions de type $(2,2,1)$ de $I$, qui \ill{} \uncertain{intersections} \ill{} les structures \uncertain{cycliques} sur quatre parmi les 5 pts de $I$
\drawing{un quadrilatère dont les quatre sommets sont marqués d'un gros point, deux côtés opposés portant chacun un petit point en leur milieu}
(\uncertain{associées} : une \uncertain{telle} \uncertain{couvrant} les deux paires \uncertain{de sommets} diagonaux \uncertain{opposés}) correspondent aux 15 arêtes dudit icosaèdre grand. \uncertain{Leurs} \add{\uncertain{a}} \uncertain{extrémités} \uncertain{contenues} \add{(\uncertain{incidentes})} \struck{dans} à une \uncertain{face}/des \ill{} l'icosaèdre, ses \og triangles\fg{} $\subset I$ correspondent : \ill{} \uncertain{commune} \uncertain{intersection} \ill{} la \uncertain{corde} \ill{}, \struck{deux est} \add{l'un des} \uncertain{sommets}, i.e. ssi le complémentaire des triangles \struck{On voit donc} \add{des deux paires \uncertain{différentes} à la corde}. \struck{\uncertain{Correspondant}} \add{\uncertain{opposés}} \ill{} \add{\uncertain{qui en} \ill{}} \uncertain{D'où} \uncertain{la relation d'incidence} \uncertain{des} \add{diviseurs} de la structure $\widehat{M}_{0,I}$ : \uncertain{correspond} : celle des faces et arêtes de l'icosaèdre grand, ou par dualité, à celle des \struck{arêtes} \add{sommets} et des \struck{faces} \add{arêtes} du dodécaèdre.

On a donc trouvé :
\marginal{NB on \uncertain{associe} à \uncertain{l'icosaèdre} \ill{} \uncertain{coloration} des faces \ill{} 4 \uncertain{couleurs}, \ill{} \uncertain{faces} \ill{} \uncertain{grand} \ill{}}
\drawing{en marge, un losange partagé par sa diagonale verticale, la moitié gauche marquée « fa »}
\note{page très raturée ; la prose entre les énoncés n'est lue qu'en partie, et l'ordre des ajouts interlinéaires reste incertain}

\page{67}\note{p. 33 de l'auteur}
\emph{Proposition.} Le diviseur à l'infini (\add{$M^{\infty}_{0,I}$}) de $\widehat{M}_{0,I}$ est une courbe \add{\uncertain{déployée}} \uncertain{stable} sur $\operatorname{Spec}\mathbb{Z}$, dont le graphe est isomorphe au \struck{celui du} \uncertain{1-squelette} du dodécaèdre grand défini par $I$ (muni d'une orientation convenable).

Bien entendu, les \uncertain{groupes} \add{des \ill{} \uncertain{liés aux} \uncertain{arêtes}} \uncertain{libres} \uncertain{aux} \uncertain{arêtes}. Il n'y a pas \add{(\uncertain{dans} \og pts marqués\fg{})} \uncertain{les} \add{\uncertain{10}} \uncertain{composantes} \ill{} \uncertain{de} \struck{$M_{0,5}$} $M^{\infty}_{0,5}$, qui \uncertain{soit} \struck{\ill{}} \add{\ill{}} \uncertain{courbe} \uncertain{de toutes façons} une \uncertain{MD}- \ill{} \uncertain{MD-graphe} \uncertain{maximal}, \uncertain{tels} \uncertain{toutes} les 2-\uncertain{sommets} \ill{} (\uncertain{dans} les gerbes \ill{} \uncertain{d'ordre 3}, \uncertain{par} \uncertain{des} \uncertain{pointes} \uncertain{marquées}).

L'isom. de la prop. \uncertain{précédente} \uncertain{n'est} \ill{} \uncertain{compatible} \uncertain{qu'avec} $\mathfrak{S}^{+}_{I}$ ${\simeq}$ \uncertain{Aut du dodécaèdre grand}, (sur $M^{\infty}_{0,I}$), mais $\mathfrak{S}_I$ \struck{tout} \uncertain{entier} -- \uncertain{des} \uncertain{opérations} de \ill{} par celles de \ill{} -- il faudrait le comprendre du pt de vue de la géométrie de l'icosaèdre (\add{grand}) \uncertain{ou} \ill{} dodécaèdre (\add{grand}) ou … -- ça doit être lié bien sûr au \emph{bi-icosaèdre} grand (\uncertain{combinatoire}).

Quand on passe à la multiplicité modulaire
\marginal{\emph{NB} $\widehat{M}_{0,I}$ \uncertain{est} un vrai schéma, \uncertain{projectif} et \uncertain{lisse} \uncertain{sur} $\mathbb{Z}$, \ill{} \ill{} \uncertain{Deligne} \ill{} \ill{} $\widehat{M}_{0,I} \simeq \widehat{M}_{0,5}$ \ill{} $V \times \mathbb{N}$…}

\page{68}
\[
(82)\qquad \widehat{M}_{0,5} = (\widehat{M}^{!}_{0,5}, \mathfrak{S}_5) = (\widehat{M}_{0,I}, \mathfrak{S}_I)
\]
il n'y a plus, par \og passage au quotient\fg{} par $\mathfrak{S}_5$, qu'un seul diviseur à l'infini, et un seul pt \add{critique} à l'infini (les 10 diviseurs à l'infini, ainsi que les 15 pts critiques à l'infini, de $\widehat{M}_{0,I}$ sont permutés transitivement entre eux par $\mathfrak{S}_I$).

\page{69}\note{p. 34 de l'auteur}
\textbf{4°) \emph{Cas de} $\widehat{M}^{!}_{1,2} \simeq \widehat{M}^{\xi}_{1,I}$} (card $I = 2$, $I = \{a, b\}$)
\note{l'exposant de $\widehat{M}_{1,I}$ est surchargé, lu $\xi$ sans certitude ; dans « card $I = 2$ » un chiffre est raturé}
\[
(83)\qquad
\begin{array}{c|c|c}
\text{codim } 0 & \text{codim } 1 & \text{codim } 2 \\ \hline
\underset{1}{\overset{I(2)}{\bullet}} &
\begin{array}{l} G_0 = \text{(boucle sur un sommet } I(2)\text{)} \\ G_1 = \overset{1}{\bullet}\!\text{---}\!\overset{I(2)}{\bullet} \end{array} &
\begin{array}{l} \overset{\{a\}}{\bullet}\!=\!\overset{\{b\}}{\bullet} = P_0 \\ \text{(boucle)}\!\bullet\!\text{---}\!\overset{I(2)}{\bullet} = P_1 \end{array}
\end{array}
\]
\drawing{tableau à trois colonnes. Codimension 1 : $G_0$ est une boucle fermée sur un sommet épais marqué $I(2)$ ; $G_1$ un segment joignant un sommet marqué 1 à un sommet marqué $I(2)$. Codimension 2 : $P_0$, deux sommets marqués $\{a\}$ et $\{b\}$ joints par deux arêtes ; $P_1$, une boucle portée par un sommet relié par une arête à un sommet marqué $I(2)$. Des flèches vont de $G_0$ vers $P_0$ (marquée 2) et vers $P_1$ (marquée 1), et de $G_1$ vers $P_1$ (marquée 1)}
\marginal{\og $D_{\circ}$ boucle\fg{} ; \og $D_{-}$ segment\fg{}}
Il y a deux diviseurs à l'infini, $D_{\circ}$ et $D_{-}$, \add{\uncertain{cf.} $P_0$}, et deux points \add{critiques} à l'infini, \struck{\ill{}} $P_0$ et $P_1$, \add{\uncertain{dont} l'un $P_0$} contenu dans $D_{\circ}$ et dans $D_{-}$, et l'autre contenu \add{\uncertain{bien}} que dans $D_{-}$.
\note{la phrase est surchargée d'ajouts ; telle qu'écrite elle place l'un des points dans les deux diviseurs et l'autre dans un seul, sans qu'on puisse dire sûrement lequel}

Essayons de \uncertain{déterminer} la \uncertain{structure de} la multiplicité \uncertain{à l'infini} $\widehat{M}^{\infty}_{1,I} \simeq \widehat{M}^{!\,\infty}_{1,2}$, \add{\uncertain{donc}} \uncertain{celle} des deux multiplicités \uncertain{\ill{}}, \struck{\ill{}} $D_{\circ}$ et $D_{-}$.

$D_{\circ}$ est isomorphe \add{MD-} \uncertain{à} $(\widehat{M}_{G_0}, \Gamma^{!}_0)$, où $G_0$ est le graphe $\bigcirc\!\bullet^{I(2)}$, $\Gamma_0$ \uncertain{le} groupe d'automorphismes, \uncertain{et} $\Gamma^{!}_0 = \operatorname{Ker}(\Gamma_0 \to \mathfrak{S}_{I(2)}) \simeq \mathbb{Z}/2$. D'ailleurs $\widehat{M}_{G_0} \simeq \widehat{M}^{!}_{0,4} \simeq \mathbb{P}^1_{\mathbb{Z}}$, sur lequel \uncertain{dès lors} $\Gamma^{!}_0 \simeq \mathbb{Z}/2$ opère en fixant un des \uncertain{pts} à l'infini \add{soit $1$} de $U_{0,3}\mathbb{Z}$ \add{\uncertain{soit} $\{0, \infty\}$}, et \uncertain{en} \uncertain{\ill{}} \uncertain{les deux autres}, \struck{(\uncertain{ceux correspondant} plus \ill{})} \struck{à $D_{-}$} \ill{} \uncertain{sur} \ill{} \uncertain{choix} d'un \uncertain{élément} de $I$ \uncertain{comme} $1$, et les deux éléments de $\widetilde{A}(G_0)$ \uncertain{comme} $0$ et $\infty$), donc \uncertain{on trouve} en \og quotient\fg{} par $z \mapsto z^{-1}$. \uncertain{Désignons}
\note{la fin de la page est très surchargée ; la ligne biffée commence par $D_{-}$}

\page{70}
plus intrinsèque, le graphe $G_0$ \struck{\uncertain{déployé}} revient à un $I \amalg \widetilde{A}^{!}_0$ de cardinal $4$, muni d'une partition de type $(2,2)$ (et $=$ \uncertain{d'une} partie $I$ de card $2$), l'une des \add{$J$} partitions du type $(1,3)$ a dans un pt marqué $\varepsilon_1$, les deux autres pts correspondant \uncertain{aux} éléments de l'\uncertain{ensemble} $J' = I \wedge_{\mathbb{Z}/2} \widetilde{A}^{!}_0$ :
\[
(84)\qquad J = \{\varepsilon_1\} \amalg J' \qquad J' = I \wedge_{\mathbb{Z}/2} \widetilde{A}^{!}_0 .
\]
Ceci posé, $\mathbb{P}^1_J = \mathbb{P}^1_{\{\varepsilon_1\} \amalg J'}$ \struck{\ill{}} peut être considéré comme le compactifié de $\mathfrak{S}^{1}_{\mathbb{Z}} \wedge_{\mathbb{Z}/2} J'$ \uncertain{(?)}, et $\mathfrak{S}_{J'} \simeq \mathbb{Z}/2$ y opère par transport de structure. D'où
\marginal{$\varepsilon_1$ correspondant au point $1$}
\[
(85)\qquad
\begin{aligned}
\widehat{M}_{G_0} &\simeq \bigl(\mathbb{P}^1_{\{\varepsilon_1\} \amalg J'}\text{\struck{$, \mathfrak{S}_{J'}$}}\bigr) \\
\underbrace{(\widehat{M}_{G_0}, \Gamma^{!}_0)}_{\widehat{M}_{D_{\circ}}} &\simeq \bigl(\mathbb{P}^1_{(\varepsilon_1 \amalg J')}, \mathfrak{S}_{J'}\bigr)
\end{aligned}
\]
\note{le symbole $\mathfrak{S}^{1}_{\mathbb{Z}}$ est une lecture douteuse}
sur lesquels $\mathfrak{S}_I$ opère via $\mathfrak{S}_{J'}$, donc \emph{$\mathfrak{S}_I$ opère trivialement sur} $\widehat{M}_{D_{\circ}}$. Les deux pts à l'infini de $\widehat{M}_{D_{\circ}}$ correspondent : la \uncertain{conjugaison} $J'$ d'une paire (pt. $P_0$) : la \uncertain{conjugaison} $\varepsilon_1$ d'autre part (point $P_1$). On a
\[
(86)\qquad M_{D_{\circ}} \simeq \mathbb{E}^1_{\mathbb{Z}} \setminus \{\text{section } 2\} = \operatorname{Spec} \mathbb{Z}[T]_{(T-2)}
\]
\[
\text{via}\quad
\begin{cases}
M_{G_0} \longrightarrow M_{D_{\circ}} \\
T \longmapsto T + \frac{1}{T}
\end{cases}
\]
\[
(87)\qquad \widehat{M}_{D_{\circ}} \simeq (\mathbb{P}^1_{\mathbb{Z}}, \text{avec singularité } 2 \text{ le long de la section } T = 2)
\]
\note{la formule (87) se poursuit à la page suivante}

\page{71}\note{p. 35 de l'auteur}
dans ce dernier iso, $P_0$ correspond à la section à l'infini, et $P_1$ à la section $2$ de $\mathbb{P}^1_{\mathbb{Z}}$.

D'autre part,
\[
(88)\qquad \widehat{M}_{D_1} \simeq \widehat{M}_{G_1} \simeq \widehat{M}_{1,1} \simeq (\mathbb{P}^1_{\mathbb{Z}}, \ldots)
\]
\note{dans (88), après la virgule, un passage biffé : « avec singularités $4$ et $6$ le long des sections $0$ et $J_1 = 2^{\alpha}3^{\beta}$ (correspond. à la courbe elliptique d'automorphisme $\mathbb{Z}/4\mathbb{Z}$) » ; la fin de (88) est barrée de grands traits obliques, ainsi qu'une suite raturée (« … des sections $\infty$, … $\widehat{M}_{D_1}$ … »), illisible ; le $D_1$ de cette page est sans doute le diviseur noté $D_{-}$ p. 69}

J'y renonce, compte tenu qu'il me faudrait d'abord comprendre la structure de $M_{1,1}$ (\add{\uncertain{c'est}} $=$ de $M_{1,1}$), \uncertain{ce que je ne connais} dans la catégorie transcendante (\uncertain{analytique} complexe, \uncertain{avec} \uncertain{éventuellement} la structure réelle explicitée -- i.e. considération de la \uncertain{conjugaison} complexe…). J'y reviendrai plus loin.

\emph{Notons déjà}

\emph{Proposition.} Si $\dim M_G \geq 2$, alors $M^{\infty}_{G}$ est connexe \add{(fibre par fibre)} \add{\uncertain{géom.}}.

\emph{Dém.} Écrivons
\[
(89)\qquad
\begin{cases}
\widehat{M}_G \simeq \prod\limits_{\alpha \in S} \widehat{M}^{!}_{g_\alpha, \hat{\nu}_\alpha} \\
M_G \simeq \prod\limits_{\alpha \in S} M^{!}_{g_\alpha, \nu_\alpha}
\end{cases}
\]
\note{sous le second produit, un indice raturé}

\page{72}
et utilisons le fait que les $\widehat{M}^{!}_{g_\alpha, \hat{\nu}_\alpha}$ sont connexes, ainsi que $\widetilde{M}^{\infty}_{G} = \widehat{M}_G \setminus M_G$ est connexe dès que \struck{deux} parmi les $\widehat{M}^{!}_{g_\alpha, \hat{\nu}_\alpha}$, il y en a au moins deux qui sont de dim $\geq 1$ (on utilise le fait que pour un tel $\alpha$, \ill{}
$M^{\infty}_{g_\alpha, \hat{\nu}_\alpha} \neq \emptyset$). Dans le cas \uncertain{contraire}, on a
\[
\widehat{M}_G \simeq \widehat{M}^{!}_{g,\nu} \quad (\simeq M_{g,I} \ldots)
\]
\marginal{$(\dim M_{g,\nu} = N = 3(g-1) + \nu \geq 2.)$}
le cas où $3(g-1) + \nu = 2$ est déjà vu, \struck{\ill{}} dans DPS $N \geq 3$. \struck{Dans ce cas}, il suffit de voir que deux diviseurs à l'infini peuvent être joints par une \uncertain{chaîne} de tels div. dont deux consécutifs ont une intersection non vide, compte tenu que ces diviseurs sont eux-mêmes connexes. Or les diviseurs à l'infini correspondent aux \struck{1-types} \add{MD-}graphes suivants
\[
(90)\quad
\begin{cases}
\text{a)}\ \ \text{une boucle sur un sommet } \overset{I}{\underset{g-1}{\bullet}} \quad (\text{si } g \geq 1) \\
\text{b)}\ \ \overset{I'}{\underset{g'}{\bullet}}\!\text{---}\!\overset{I''}{\underset{g''}{\bullet}} \quad
\begin{cases}
I = I' \amalg I'' \\
g' + g'' = 0 \\
\text{si } g' \text{ (ou } g'') = 0, \text{ alors} \\
\quad \text{\struck{$2g' + \operatorname{card} I' \geq 2$}} \\
\quad \text{\struck{$2g'' + \operatorname{card} I'' \geq 2$}} \\
\quad \operatorname{card} I' \geq 2 \ (\text{ou } \operatorname{card} I'' \geq 2)
\end{cases}
\end{cases}
\]
\note{dans b), un premier dessin raturé à gauche du segment ; « $g' + g'' = 0$ » se lit sur la page, sans doute pour $g$}

\page{73}\note{p. 36 de l'auteur}
On a comme spécialisation des cas a)
\[
(91)\quad \text{a}')\quad
\text{une boucle sur } \overset{I_1}{\underset{g_1}{\bullet}}\!\text{---}\!\overset{I_2}{\underset{g_2}{\bullet}}
\qquad
\begin{cases}
I_1 \amalg I_2 = I \\
g_1 + g_2 = g - 1 \\
\text{si } g_2 = 0 \text{ \struck{\ill{}}}\ \text{\add{on a card } } I_2 \geq 2 \\
\quad \text{\struck{($I_1 \neq \emptyset$ et card $I_2 \neq \emptyset$)}}
\end{cases}
\]
\note{au-dessus de l'arête du dessin, un trait ondulé raturé}
\struck{\ill{}} DPS $g_2 = 1$, $g_1 = g - 2$, $(I_1, I_2)$ partition quelc. de $I$ (pourvu que $g \geq 2$).
On a comme \uncertain{autres} spécialisation des cas b), \struck{si \ill{} \ill{} $g'$}
\[
(92)\quad \text{b}')\quad
\text{une boucle sur } \overset{I'}{\underset{g'-1}{\bullet}}\!\text{---}\!\overset{I''}{\underset{g''}{\bullet}}
\]
Donc si $g \geq 1$ (\uncertain{donc} \uncertain{dès} que le cas a) existe), dans presque tous les cas b) \struck{$g' \neq$} les \uncertain{conditions} \ill{} \uncertain{de} $g', g''$ ne sont que $g' \neq 0$ i.e. $g' \geq 1$ (quitte à éch. $g'$, $g''$), on voit que l'on a \add{$b') = a')$} une spécialisation commune aux cas a) et b). Cela démontre notre assertion lorsque $g \neq 0$. \struck{Si $g$} Dans le cas $g = 0$, \add{\uncertain{donc} card $I \geq 6$}, \struck{\uncertain{de} DPS} on a cas les cas b), \struck{a il y a} \add{\uncertain{dès qu'}il y a} \struck{que} card $I' \geq$ card $I''$ \struck{dans} card $I' \geq 3$, une spécialisation de la forme
\[
\text{b}'')\qquad
\text{\struck{(b'') : \ill{}}}\quad
\overset{\mathfrak{f}}{\bullet}\!\text{---}\!\bullet
\]
\drawing{schéma raturé : un graphe à deux sommets marqués $I' \setminus \{a'\}$ et $I''$, entouré d'un contour en forme de cloche ; en dessous, un segment dont l'extrémité gauche porte un symbole surchargé}
\[
(93)\quad \text{b}'')\quad
\underbrace{\overset{2}{\underset{(a_0)}{\bullet}}\!-\!\overset{1}{\underset{(a_1)}{\bullet}}\!-\!\overset{1}{\bullet}\cdots\overset{1}{\underset{(a_{k-1})}{\bullet}}\!-\!\overset{1}{\underset{a_k}{\bullet}}}_{\substack{k \text{ segments},\ k+1 \text{ sommets} \\ \text{ayant un poids marqué total} \\ 2 + k = \operatorname{card} I'}}
\!-\!
\underbrace{\overset{1}{\underset{(b_l)}{\bullet}}\!-\!\overset{1}{\underset{(b_{l-1})}{\bullet}}\cdots\overset{1}{\underset{(b_2)}{\bullet}}\!-\!\overset{1}{\underset{(b_1)}{\bullet}}\!-\!\overset{2}{\underset{(b_0)}{\bullet}}}_{\substack{l \text{ segments},\ l+1 \text{ sommets} \\ \text{poids marqué total} \\ 2 + l = \operatorname{card} I''}}
\]
\note{sous le second groupe on lit « $2 + l = \operatorname{card} I''$ », le premier chiffre surchargé ; les poids au-dessus des sommets sont les nombres de points marqués}

\page{74}
diagrammes qui ont \uncertain{ceci} de \uncertain{commun}, \uncertain{à savoir} qu'ils \uncertain{sont} \uncertain{donnés} par un ordre total sur $I$, \uncertain{(défini} \uncertain{à} \uncertain{transposition} \struck{\ill{}} \uncertain{près} des deux premiers et des deux derniers). Pour montrer que deux ordres totaux sur $I$ définissent \struck{\ill{}} des pts de $M^{\infty}_{0,\nu}$ qui \uncertain{appartiennent} à une \uncertain{même} composante connexe, il suffit de \uncertain{montrer} que \uncertain{cela} \uncertain{est vrai} pour deux ordres totaux \uncertain{déduits} l'un de l'autre \uncertain{par} \uncertain{transposition} des deux éléments $c_i, c_{i+1}$ consécutifs \uncertain{dans} \uncertain{les} \uncertain{deux} $=$ \uncertain{composante} connexe. Or les \uncertain{deux} \uncertain{sont} \uncertain{toujours} \uncertain{la} \uncertain{spécialisation} de
\[
\overset{(c_1, c_2)}{\underset{c}{\bullet}}\!-\!\overset{c_3}{\bullet}\!-\!\cdots\!-\!\overset{c_{i-1}}{\bullet}\!-\!\overset{(c_i, c_{i+1})}{\bullet}\!-\!\overset{c_{i+2}}{\bullet}\!-\!\cdots\!-\!\overset{c_\nu}{\bullet}
\]
\note{devant cette chaîne, un mot raturé ; on lit « dans »}
provenant de
\[
\cdots\!-\!\overset{c_{i-1}}{\bullet}\!-\!\overset{c_i}{\bullet}\!-\!\overset{c_{i+1}}{\bullet}\!-\!\overset{c_{i+2}}{\bullet}\!-\!\cdots
\]
et de
\[
\cdots\!-\!\overset{c_{i-1}}{\bullet}\!-\!\overset{c_{i+1}}{\bullet}\!-\!\overset{c_i}{\bullet}\!-\!\overset{c_{i+2}}{\bullet}\!-\!\cdots
\]
en contractant la \struck{\ill{}} \add{arête} qui joint $c_i$ et $c_{i+1}$, OK.

\page{75}\note{p. 37 de l'auteur}
\textbf{9) \emph{La structure groupoïdale des Teichmüller modulaires} (\add{\uncertain{groupes} \uncertain{modulaires} \uncertain{des}}) \emph{(la MDT-structure) : cas transcendant}}
\note{l'insertion au-dessus du titre est soulignée et en partie illisible}
(morphismes de généralisation.)

Nous travaillons maintenant avec les $M^{\mathrm{an}}_{g,\nu}$, $\widehat{M}^{\mathrm{an}}_{g,\nu}$ etc. -- et pour ne pas traîner les exposants \og an\fg{} (pour \og analytique\fg{}), nous les laisserons tomber.

Considérons pour tout MD-graphe $\underline{G}$, la multiplicité modulaire (analytique) $M_{\underline{G}}$, et son groupoïde fondamental, que je note $\Pi M_{\underline{G}}$. C'est une \og \add{\uncertain{grosse}}\fg{} \struck{catégorie}, par le fait même \uncertain{qu'elle} est \uncertain{connexe} \uncertain{gouverne} le \uncertain{remplacer} par un \struck{catégorie} \add{groupoïde} \uncertain{dont} l'ensemble des objets soit fini. Dans nos constructions, il est en fait licite de remplacer $\Pi M_{\underline{G}}$ par un groupoïde tel que \struck{l'équivalence} \add{des \uncertain{échantillons}} groupoïde équivalent, pourvu que \struck{\ill{}} l'équivalence commute à l'action \uncertain{du groupe}
\[
(94)\qquad \Gamma_{\underline{G}} = \operatorname{Aut} \underline{G}
\]
\note{devant (94), une ligne « le groupe » rayée}

Notons déjà que $\Pi M_{\underline{G}}$ est un groupoïde
\note{la phrase se poursuit à la page suivante}

\page{76}
connexe, dont le $\pi_1$ est isomorphe au groupe \emph{(\uncertain{cartésien})} :
\[
(95)\qquad \pi_1(\Pi M_{\underline{G}}) = \pi_1(M_{\underline{G}}) \simeq \prod_{\alpha \in S(\underline{G})} \pi_1(M_{g_\alpha, \hat{I}_\alpha}) \simeq \prod_\alpha \Gamma^{+}_{g_\alpha, I_\alpha} \simeq \prod_\alpha \mathcal{T}^{!+}_{g_\alpha, \nu_\alpha}
\]
\marginal{les $\mathcal{T}^{!+}_{g,\nu}$ sont \uncertain{dits} \uncertain{des} \uncertain{groupes} de Teichmüller}
\note{le dernier isomorphisme est écrit sous le précédent, relié par un signe $\simeq$ vertical ; l'indice $\hat{I}_\alpha$ est surchargé}
\uncertain{L'opération} de $\Gamma^{\natural}_{\underline{G}}$ sur $\Pi M_{\underline{G}}$ définit un groupoïde \uncertain{muni} $(\Pi M_{\underline{G}}, \Gamma_{\underline{G}})$, dont le groupe fondamental (\uncertain{cartésien}) \struck{\ill{}} \add{\uncertain{dans} \uncertain{une}} structure d'extension
\[
(96)\qquad 1 \to \pi_1(\Pi M_{\underline{G}}, \Gamma_{\underline{G}}) \to \pi_1(\Pi M_{\underline{G}}, \Gamma_{\underline{G}}) \longrightarrow \Gamma_{\underline{G}} \to 1
\]
\note{tel qu'écrit, le premier terme non trivial est le même que le second ; on attend $\pi_1(\Pi M_{\underline{G}})$}

Nous allons définir une \add{\uncertain{sorte de}} \uncertain{variation} des $\Pi M_{\underline{G}}$ relativement aux \og morphismes\fg{} pour les MD-graphes $\underline{G}$. Il s'agit d'abord \uncertain{d'une} \og \emph{variation de généralisation}\fg{} \uncertain{qui} est comprise de la \emph{variation pour morphismes des MD-graphes} (que j'admets comme \og \uncertain{triviale}\fg{}, et qu'il est inutile de préciser) et une \og variation de \struck{\ill{}} \uncertain{généralisation} de MD-graphes\fg{}.
\note{la fin de la page est d'une lecture très incertaine : le texte distingue une variation par les morphismes de MD-graphes et une variation de généralisation}

\page{77}\note{p. 38 de l'auteur}
Soit donc $\underline{G}'$ un MD-graphe quotient de $\underline{G}$ -- par un ens. $C$ d'arêtes du graphe $G$ sous-jacent à $\underline{G}$,
\[
(96)\qquad \underline{G}' = \underline{G} / C .
\]
\note{le numéro se lit (96), déjà employé p. 76}
On a alors un \add{diagramme de} morphismes canoniques

\begin{tikzcd}
M_{\underline{G}} \arrow[d, hook] \\
\widehat{M}_{\underline{G}} \arrow[r] & \widehat{M}_{\underline{G}'}
\end{tikzcd}

\note{(97) : à droite de $M_{\underline{G}}$, une flèche vers $M_{\underline{G}'}$ et une flèche verticale de $M_{\underline{G}'}$ vers $\widehat{M}_{\underline{G}'}$ sont biffées de grands traits}
\marginal{diptyque et gît en $\Gamma^{!}_{\underline{G},\underline{G}'}$}
\struck{\ill{}} \uncertain{Si l'on} désigne par $\Gamma^{!}_{\underline{G},C} = \Gamma^{!}_{\underline{G},\underline{G}'}$ le sous-groupe de $\Gamma_{\underline{G}}$ qui est formé des \uncertain{automorphismes} \uncertain{qui} \uncertain{induisent} l'identité sur $\underline{G}'$, \uncertain{alors} le \uncertain{morphisme} \struck{\ill{}} \add{horizontal} de (97) se factorise \uncertain{en} \uncertain{une} \uncertain{immersion} de multiplicités
\[
(98)\qquad (\widehat{M}_{\underline{G}}, \Gamma^{!}_{\underline{G},\underline{G}'}) \hookrightarrow \widehat{M}_{\underline{G}'}
\]
\marginal{que \uncertain{ce} \uncertain{ne} \uncertain{sont} \uncertain{peut-être} \uncertain{que} \uncertain{des} \uncertain{immersions} \uncertain{locales}, il \uncertain{paraît} \ill{}…}
Je désigne par $\widehat{M}_{\underline{G},\underline{G}'}$ le \emph{germe de} multiplicité défini par l'inclusion (98), et par $M_{\underline{G},\underline{G}'}$ le germe défini par l'inclusion
\[
(99)\qquad (M_{\underline{G}}, \Gamma^{!}_{\underline{G},\underline{G}'}) \hookrightarrow \widehat{M}_{\underline{G}'}
\]
\marginal{Il vaut mieux considérer le \uncertain{germe} de multiplicité, \uncertain{à} \uncertain{définir} \uncertain{par} $M_{\underline{G}} \to M_{\underline{G}'}$ \uncertain{dessus} $\Gamma^{!}_{\underline{G},\underline{G}'}$ \ill{}}
\note{deux notes marginales obliques, à gauche, sont d'une lecture très incertaine}

\page{78}
Je désigne par $\widehat{M}^{*}_{\underline{G},\underline{G}'}$, $\widehat{M}^{*}_{\underline{G},\underline{G}'}$ les germes de multiplicités obtenus en \struck{prenant les traces des $\widehat{M}_{\underline{G}'}$ \ill{}} \add{enlevant les traces des diviseurs \uncertain{à l'infini} de $\widehat{M}_{\underline{G}'}$} \struck{\og l'infini\fg{} \ill{}} \struck{les traces, \ill{} \ill{} de $M_{\underline{G}}$ \ill{}} \ill{} \uncertain{filtrées} \uncertain{d'}\uncertain{éléments} de la multiplicité $M_{\underline{G}'}$, et les groupoïdes fondamentaux* \uncertain{\ill{}}
\note{le second symbole porte un accent surchargé, peut-être un double chapeau ; tout le haut de la page est couvert de ratures et d'ajouts entrecroisés, dont l'ordre reste douteux}
\marginal{\uncertain{qui} \uncertain{sont} \uncertain{\ill{}} \uncertain{dans} le \uncertain{cas} \uncertain{présent} les $M^{*}_{\underline{G},\underline{G}'}$ $M_{\underline{G}'}$ \uncertain{dans} \uncertain{\ill{}} $\widehat{M}_{\underline{G},\underline{G}'}$}
\[
(100)\qquad \Pi M^{*}_{\underline{G},\underline{G}'}, \quad \Pi M^{*}_{\underline{G},\underline{G}'}
\]
\note{les deux termes de (100) se lisent identiques ; sans doute l'un porte-t-il un chapeau}
sont définis par \uncertain{une} \uncertain{des} catégories $\varinjlim$ des catégories de revêtements \uncertain{universels}, ou des catégories des \uncertain{revêtements} \uncertain{connexes} desdits germes de multiplicités (\uncertain{comparer} [53], n°  5, p. 768').

On a alors des diagrammes de germes de multiplicités \uncertain{analytiques}
\[
(101)\qquad
\begin{array}{ccc}
M^{*}_{\underline{G},\underline{G}'} & \longrightarrow & M_{\underline{G}'} \\
\big\downarrow & & \\
(M_{\underline{G}}) \overset{\text{\struck{\ill{}}}}{\longrightarrow} M_{\underline{G},\underline{G}'} & &
\end{array}
\]
\note{sous la flèche horizontale du bas : « (équivalence d'homotopie !) » ; un indice est raturé après $M_{\underline{G}'}$}
donnant lieu : \uncertain{des} morphismes de groupoïdes
\note{la phrase se poursuit à la page suivante}

\page{79}\note{p. 39 de l'auteur}
\[
(102)\qquad
\begin{array}{ccc}
\Pi M^{*}_{\underline{G},\underline{G}'} & \longrightarrow & \Pi M_{\underline{G}'} \\
\big\downarrow & & \\
\Pi(M_{\underline{G}}) \overset{\approx}{\longrightarrow} \Pi M_{\underline{G},\underline{G}'} & &
\end{array}
\]
\note{sous la flèche du bas, une note de sa main : « ceci est une équivalence de groupoïdes » ; dans le terme de gauche, un symbole est raturé}
qui se réduisent \uncertain{à} un \uncertain{morphisme} \uncertain{de} \uncertain{groupoïdes}
\[
(103)\qquad
\begin{array}{ccc}
\Pi M^{*}_{\underline{G},\underline{G}'} & \xrightarrow{\ \varphi_{\underline{G},\underline{G}'}\ } & \Pi M_{\underline{G}'} \\
\big\downarrow{\scriptstyle \varphi_{\underline{G},\underline{G}'}} & & \\
\Pi(M_{\underline{G}}, \ldots) & &
\end{array}
\]
\note{l'étiquette de la flèche horizontale est surchargée ; à droite du terme du bas, un long passage est biffé de grands traits obliques, où l'on devine « sera \ill{} … $\Pi M_{\underline{G},\underline{G}'}$ … $\Pi(M_{\underline{G}}, \Gamma^{!}_{\underline{G},\underline{G}'})$ »}
\marginal{\og \uncertain{diagramme} de \uncertain{généralisation}\fg{} pour \uncertain{les} groupoïdes \uncertain{fondamentaux} par \uncertain{la} \uncertain{généralisation} de $\underline{G}$ \uncertain{à} $\underline{G}'$}
L'homomorphisme de groupoïdes $\varphi_{\underline{G},\underline{G}'}$ \struck{\ill{}} \uncertain{a} \uncertain{comme} \uncertain{propriétés} \struck{de} l'homomorphisme sur les $\pi_1$ des \uncertain{bouts} \uncertain{de} $\Pi(M_{\underline{G}},$ \struck{…}$)$ \struck{\ill{}} \uncertain{une} fibration de \uncertain{fibre} un produit de disques épointés (\uncertain{un} \uncertain{par} \uncertain{arête} \uncertain{de} $C$, \uncertain{ou} \uncertain{plutôt}) -- on peut \uncertain{\ill{}} $\Pi M^{*}_{\underline{G},\underline{G}'}$ comme une \og extension\fg{} canonique du groupoïde \ill{} $\Pi(M_{\underline{G}}, \ldots)$ par $\mathbb{Z}^{C}$, où $C$, \uncertain{rappelons}-le, est l'ens. des \struck{flèches} \add{arêtes} de $\underline{G}$
\marginal{\emph{NB} $(M_{\underline{G}}, \ldots)$ \uncertain{est} un $K(\pi, 1)$, et $\pi_1$ \uncertain{des} \uncertain{fibres} \ill{} \uncertain{extension} \ill{} \ill{} de \uncertain{\ill{}} \uncertain{vers} \uncertain{l'}\uncertain{une}…}
\marginal{\uncertain{Notation} : \uncertain{il} \uncertain{vaut} \uncertain{mieux} \uncertain{noter} $(M_{\underline{G}}, \Gamma^{!}_{\underline{G},\underline{G}'})$ \uncertain{pour} $\Gamma^{!}_{\underline{G},\underline{G}'}$ \ill{} $\mathbb{Z}^{C}$ \ill{}}
\note{les deux notes marginales, écrites de biais et en partie biffées, sont reliées au texte par une longue flèche courbe ; leur lecture est très incertaine}

\page{80}
qui \uncertain{a été} \uncertain{introduite} pour \uncertain{traiter} la généralisation $\underline{G}'$. \struck{Plus} (On \uncertain{revient} sur la structure de cette extension, comme un produit fibré d'extensions, un peu plus \uncertain{loin}…) L'homomorphisme $\varphi_{\underline{G},\underline{G}'}$ s'appelle \emph{l'homomorphisme de généralisation} pour les \emph{groupoïdes fondamentaux}, associé à l'homomorphisme de généralisation $\underline{G} \to \underline{G}'$.

On va dégager une propriété de transitivité pour un homomorphisme de généralisation, ou \uncertain{plutôt} pour un \uncertain{composé} d'homomorphismes de généralisation
\[
(104)\qquad \underline{G} \longrightarrow \underline{G}' \longrightarrow \underline{G}''
\]
qui donne \uncertain{\ill{}} \uncertain{l'}\uncertain{association} : un diagramme de groupoïdes
\[
\text{\struck{$\Pi M^{*}_{\underline{G},\underline{G}'} \longrightarrow \ldots$ \quad $\big\downarrow$ \quad $\Pi M_{\underline{G}}$}}
\]
\note{le début de diagramme est biffé de traits obliques ; la page s'arrête là, et l'argument se poursuit au-delà de ce lot}

\end{document}
